%%=============================================================================
%% Samenvatting
%%=============================================================================

% TODO: De "abstract" of samenvatting is een kernachtige (~ 1 blz. voor een
% thesis) synthese van het document.
%
% Een goede abstract biedt een kernachtig antwoord op volgende vragen:
%
% 1. Waarover gaat de bachelorproef?
% 2. Waarom heb je er over geschreven?
% 3. Hoe heb je het onderzoek uitgevoerd?
% 4. Wat waren de resultaten? Wat blijkt uit je onderzoek?
% 5. Wat betekenen je resultaten? Wat is de relevantie voor het werkveld?
%
% Daarom bestaat een abstract uit volgende componenten:
%
% - inleiding + kaderen thema
% - probleemstelling
% - (centrale) onderzoeksvraag
% - onderzoeksdoelstelling
% - methodologie
% - resultaten (beperk tot de belangrijkste, relevant voor de onderzoeksvraag)
% - conclusies, aanbevelingen, beperkingen
%
% LET OP! Een samenvatting is GEEN voorwoord!

%%---------- Nederlandse samenvatting -----------------------------------------
%
% TODO: Als je je bachelorproef in het Engels schrijft, moet je eerst een
% Nederlandse samenvatting invoegen. Haal daarvoor onderstaande code uit
% commentaar.
% Wie zijn bachelorproef in het Nederlands schrijft, kan dit negeren, de inhoud
% wordt niet in het document ingevoegd.

\IfLanguageName{english}{%
\selectlanguage{dutch}
\chapter*{Samenvatting}
\lipsum[1-4]
\selectlanguage{english}
}{}

%%---------- Samenvatting -----------------------------------------------------
% De samenvatting in de hoofdtaal van het document

\chapter*{\IfLanguageName{dutch}{Samenvatting}{Abstract}}

Dit onderzoek gaat na in welke mate een AI-gebaseerd systeem op basis van mark\-er\-loze video-analyse technisch bruikbare feedback kan genereren voor recreatieve turnsters bij de kip aan de barre. Binnen recreatieve gymnastiekclubs begeleiden trainers vaak meerdere turnsters tegelijk, waardoor technische fouten niet altijd systematisch worden opgemerkt. Tegelijk is gerichte feedback belangrijk om de uitvoering van complexe elementen veilig en efficiënt aan te leren.

Het doel van dit onderzoek is het ontwikkelen en verkennend evalueren van een proof-of-concept dat standaard videobeelden omzet naar pose-data, daaruit relevante gewrichtshoeken berekent, vooraf gedefinieerde technische fouten regelgebaseerd detecteert en deze detecties omzet naar tekstuele feedback via een Large Language Model. De nadruk ligt daarbij niet op een volledig gevalideerd coachesysteem, maar op de technische haalbaarheid en beperkingen van een dergelijke analysepipeline.

De proof-of-concept werd opgebouwd in verschillende stappen. Eerst werd Media\-Pipe Pose Landmarker gebruikt om lichaamslandmarks uit videobeelden te extraheren. Vervolgens werden op basis van deze landmarks de heuphoek, schouderhoek, romphoek en ellebooghoek berekend en opgeslagen in CSV-formaat. Daarna werd een Python-script ontwikkeld dat een manueel gekozen analysevenster verwerkt, kritieke ankerpunten detecteert, de beweging opdeelt in fasen en vier vooraf bepaalde fouttypes analyseert: gebogen armen, onvoldoende heupsluiting, te vroege heupopening en onvoldoende schouderdruk in de steunfase. De analyse-output werd gestructureerd opgeslagen in JSON-formaat en in gefilterde vorm gebruikt als input voor een lokale LLM via Ollama.

Uit de resultaten blijkt dat markerloze pose-estimatie bruikbare bewegingsdata kan opleveren uit standaard videobeelden, maar dat de kwaliteit afhankelijk blijft van de opname. Achtergrondbewegingen, tijdelijk minder zichtbare lich\-aams\-delen en detectieruis kunnen abrupte afwijkingen veroorzaken in de berekende hoekcurves. Daarom werd de analyse iteratief verfijnd, onder meer door smoothing toe te passen en foutregels aan te passen om foutieve detecties te beperken.

De finale analyse werd getest op drie pogingen: één goede uitvoering en twee foutieve uitvoeringen. Bij de goede uitvoering werden geen fouten gedetecteerd. Bij de eerste foutieve uitvoering werden gebogen armen in de trekfase en onvoldoende schouderdruk in de steunfase vastgesteld. Bij de tweede foutieve uitvoering werd vooral onvoldoende schouderdruk in de steunfase gedetecteerd. De resultaten tonen aan dat de regelgebaseerde aanpak binnen de beperkte testset relevante verschillen tussen pogingen kan herkennen, maar niet als formele validatie mag worden beschouwd.

De LLM-feedback bleek technisch haalbaar, maar niet volledig betrouwbaar zonder strikte afbakening. Wanneer het taalmodel te veel informatie kreeg, bestond het risico dat het niet-gedetecteerde fouten of algemene verbeterpunten formuleerde. Daarom werd de LLM-input beperkt tot de effectief gedetecteerde fouten en wordt geen LLM-call uitgevoerd wanneer het regelgebaseerde systeem geen fouten detecteert.

Dit onderzoek toont aan dat automatische technische feedback op de kip aan de barre in beperkte mate haalbaar is via een combinatie van markerloze pose-estimatie, regelgebaseerde foutdetectie en LLM-gebaseerde tekstgeneratie. De aanpak biedt potentieel als ondersteunend hulpmiddel voor trainers, maar vervangt geen menselijke expertise. Verdere evaluatie op een grotere dataset en met feedback van ervaren coaches blijft noodzakelijk.
